
\chapter{Conceptual Framework}

This chapter defines the main concepts and classifications used throughout the empirical analysis. Because the thesis compares the United States and the Euro area, it is necessary to harmonize institutional boundaries, sectoral classifications, financial-institution groups, and financial instruments across the two regions. The chapter also introduces the link between national financial accounts and FWTW matrices. Together, these definitions provide the basis for constructing comparable financial exposure matrices and for attributing environmental pressures generated by real-economy agents to the financial institutions that finance them.

\section{The Euro Area}

As of 2024, the Euro area comprises 20 European Union Member States that have formally adopted the euro as their single currency: Austria, Belgium, Croatia, Cyprus, Estonia, Finland, France, Germany, Greece, Ireland, Italy, Latvia, Lithuania, Luxembourg, Malta, the Netherlands, Portugal, Slovakia, Slovenia, and Spain.\footnote{Available at: \url{https://www.ecb.europa.eu/euro/intro/html/index.en.html}} Notably, because Croatia adopted the euro on January 1, 2023, the European Central Bank's aggregate data reflects a structural expansion from 19 to 20 member states starting in the 2023 reporting periods. In this thesis, the terms \textit{Euro area} and \textit{Eurozone (EZ)} are treated as synonymous and used interchangeably.

\section{National financial accounts}

Financial accounts are an integral component of national accounts, a system of accounts and balance sheets that provide a broad and integrated framework to describe an economy of a country or a region. To guarantee cross-border comparability, these accounts rely on the global guidelines and terminology defined in the System of National Accounts (SNA 2008), with which The European System of National Accounts (ESA 2010) is fully consistent.

Financial accounts provide a comprehensive snapshot of outstanding assets and liabilities, disaggregated by institutional sector and financial instrument. By design, these accounts form a closed system where aggregate claims and obligations reconcile, ensuring that the balancing item accurately reflects the net worth of the total economy. This comprehensive coverage makes financial accounts an important tool for analyzing the financial system in its entirety. Furthermore, the framework guarantees internal consistency, thereby preventing the double counting of exposures across interconnected institutions. 

The financial accounts of the United States are compiled by the Board of Governors of the Federal Reserve System and are accessible via the quarterly \href{https://www.federalreserve.gov/releases/z1/}{Z.1 Statistical Release}. Similarly, the European Central Bank provides quarterly financial account data for both the broader Euro area economy and individual European Union (EU) Member States through its \href{https://data.ecb.europa.eu/data/publications/EAA02_02/data-information?publications_array%5B0%5D=EAA02_02&layerType=AL}{online data portal}.

\section{Real-economy Agents}
\label{Section 3.1.2}

According to the latest \textit{Financial Accounts of the United States (Z.1)} released by the Federal Reserve, \nocite{board_of_governors_of_the_federal_reserve_system_financial_2026}the US economy comprises of Domestic nonfinancial sectors, Domestic financial sectors and Rest of the World sector (RoW). Domestic nonfinancial sectors are further broken into Household and nonprofit organizations (HH), Nonfinancial business (NFC), and General government (GG). 

The RoW sector comprises all non-resident entities that engage in transactions with domestic residents. In the financial accounts, this sector is recorded from the perspective of non-residents, which allows the RoW and domestic sectors to be treated consistently as both suppliers and users of funds. Accordingly, the acquisition of domestic assets by either RoW or domestic sectors represents a source of liquidity for domestic capital markets, while increases in the liabilities and equity of these sectors reflect funding provided through domestic markets.

The ESA 2010 adopts a conceptually equivalent sectoral structure, with minor differences in terminology. To ensure comparability, I group non-financial sectors into four harmonized agent categories (HH, NFC, GG, and RoW), indexed by $k = 1, \dots, n_R$. The correspondence in terminology across the US and Euro area classifications is summarized in Table \ref{tab:sector_mapping}. 

\begin{table}[ht]
\centering
\caption{Real-economy Agents: US and Eurozone}
\label{tab:sector_mapping}
\small
\begin{tabularx}{\textwidth}{@{} c >{\raggedright\arraybackslash}X >{\raggedright\arraybackslash}X @{}}
\toprule
\textbf{Agent Group} & \textbf{United States} & \textbf{Eurozone} \\
\midrule
% ---------------------------------------------------------
% 1. HOUSEHOLDS (HH)
% ---------------------------------------------------------
\textbf{HH} & Households and nonprofit organizations & Households \\
 & &Non-profit institutions serving households \\
\midrule
% ---------------------------------------------------------
% 2. NON-FINANCIAL CORPORATIONS (NFC)
% ---------------------------------------------------------
\textbf{NFC} & Nonfinancial business & Non-financial corporations \\
 & \quad \textit{Nonfinancial corporate business} & \\
 & \quad \textit{Nonfinancial noncorporate business} & \\
\midrule
% ---------------------------------------------------------
% 3. GENERAL GOVERNMENT (GG)
% ---------------------------------------------------------
\textbf{GG} & General government & General government \\
 & \quad \textit{Federal government} & \quad \textit{Central government} \\
 & \quad \textit{State and local governments} & \quad \textit{State government} \\
 & & \quad \textit{Local government} \\
 & & \quad \textit{Social security funds} \\
\midrule
% ---------------------------------------------------------
% 4. REST OF THE WORLD (RoW)
% ---------------------------------------------------------
\textbf{RoW} & Rest of the world & Rest of the world \\
\bottomrule
\multicolumn{3}{@{}l}{\footnotesize \textit{Source: Financial Accounts of the United States (Z.1) and the European System of Accounts (ESA 2010)}}
\end{tabularx}
\end{table}

\section{Financial Institutions}
\label{sec:fi_def}
The classification of financial institutions is more diverse and reflects fundamental differences in financial market structures across the two regions. As shown in Table \ref{tab:fi_mapping}, the categories of financial institutions in the United States are more detailed and complex than those in the Euro area. Therefore, my approach is to start with grouping the data according to the Euro area classification and then aggregate the US financial institutions accordingly.

\textbf{Monetary financial institutions}
as defined by the ECB comprise the central bank and \textit{"financial intermediaries through which the effects of the monetary policy of the central bank are transmitted to the other entities of the economy"} (Article 2.68, ESA 2010). They are deposit-taking corporations and money market funds. The latter are classified within monetary financial institutions because their primary activity is issuing investment fund shares that are close substitutes for deposits, while investing mainly in money market instruments or short-term debt securities. In the United States, comparable institutions can be identified, with the central bank classified under the "Monetary Authority" sector, represented by the Federal Reserve System. 

\textbf{Non-money market investment funds} issue shares that are not close substitutes for deposits and allocate their portfolios primarily to financial assets other than short-term ones or non-financial assets. ESA 2010 identifies examples such as open- and closed-ended funds, real estate investment funds, hedge funds, and “funds of funds.” This classification allows us to group US entities, such as mutual funds (an open-ended form), closed-end funds, exchange-traded funds, and mortgage real estate investment trusts, within this category, as they closely align with the ESA definition in terms of both asset allocation and the nature of the instruments they issue.

\textbf{Insurance corporations} and \textbf{pension funds} naturally constitute two distinct categories. In the United States, property–casualty and life insurance companies are grouped under insurance corporations, while both private and public pension funds are classified under pension funds to reduce complexity in the analysis.

The remaining financial institutions not classified elsewhere are grouped under \textbf{Other financial institutions} in both regions. While US components are relatively self-explanatory, the Euro area classifications are less granular.  Under the ESA 2010 framework, this group comprises three main sub-sectors: "Other financial intermediaries" (which include securitization vehicles, securities and derivatives dealers, lending institutions, and specialized finance companies such as venture capital firms); "Financial auxiliaries" (covering brokers, payment institutions, and exchanges); and "Captive financial institutions and money lenders" (comprising holding companies, sovereign wealth funds, and special purpose entities that raise funds for their parent companies). Although the exact institutional boundaries are not perfectly aligned between the two jurisdictions, the economic functions and activities of these entities are broadly analogous. Consequently, consolidating them into a single OFI category is a methodologically sound approach that enables robust cross-regional comparison.

In summary, there are five categories of financial corporations: Monetary financial institutions (MFI), Non-money market funds investment funds (NM\_MFIF), Insurance corporations (IS), Pension funds (PEN) and Other financial institutions (OFI), indexed by $f = 1, \dots, n_F$.  

\begin{table}[ht!]
\centering
\caption{Financial Institutions: US and Eurozone}
\label{tab:fi_mapping}
\small
\begin{tabular}{@{} >{\raggedright\arraybackslash}p{3.5cm} >{\raggedright\arraybackslash}p{6cm} >{\raggedright\arraybackslash}p{5.5cm} @{}}
\toprule
\textbf{Financial Institution Group} & \textbf{United States} & \textbf{Eurozone} \\
\midrule
% ---------------------------------------------------------
% 1. MONETARY FINANCIAL INSTITUTIONS (MFI)
% ---------------------------------------------------------
\textbf{Monetary financial institutions (MFI)} & 
Monetary Authority \newline 
Private depository institutions \newline
\hspace*{1em} \textit{U.S.-Chartered Depository} \newline 
\hspace*{1em} \textit{Institutions} \newline 
\hspace*{1em} \textit{Foreign Banking Offices in U.S.} \newline 
\hspace*{1em} \textit{Banks in U.S.-Affiliated Areas} \newline 
\hspace*{1em} \textit{Credit Unions} \newline 
Money Market Funds & 
Central bank \newline 
Monetary financial institutions other than central banks \newline
\hspace*{1em} \textit{Deposit-taking
corporations}\newline
\hspace*{1em} \textit{except the
central bank}\newline
\hspace*{1em} \textit{Money Market Funds} \\

\midrule
% ---------------------------------------------------------
% 2. NON-MMF INVESTMENT FUNDS (NM_MFIF)
% ---------------------------------------------------------
\textbf{Non-money market funds investment funds (NM\_MFIF)} & 
Mutual Funds \newline 
Closed-End Funds \newline 
Exchange-Traded Funds \newline 
Mortgage Real Estate Investment Trusts & 
Non-money market funds investment funds \\
\midrule
% ---------------------------------------------------------
% 4. INSURANCE CORPORATIONS (IS)
% ---------------------------------------------------------
\textbf{Insurance corporations (IS)} & 
Property-Casualty Insurance Companies \newline 
Life Insurance Companies & 
Insurance corporations \\
\midrule
% ---------------------------------------------------------
% 5. PENSION FUNDS (PEN)
% ---------------------------------------------------------
\textbf{Pension funds (PEN)} & 
Private and Public Pension Funds & 
Pension funds \\
\midrule
% ---------------------------------------------------------
% 3. OTHER FINANCIAL INSTITUTIONS (OFI)
% ---------------------------------------------------------
\textbf{Other financial institutions (OFI)} & 
Government-Sponsored Enterprises \newline 
Agency- and Government-sponsored Enterprise (GSE)-Backed Mortgage Pools \newline 
Issuers of Asset-Backed Securities \newline 
Finance Companies \newline 
Security Brokers and Dealers \newline 
Holding Companies \newline 
Other Financial Business & 
Other financial intermediaries, except insurance corporations and pension funds \newline 
Financial auxiliaries \newline 
Captive financial institutions and money lenders \\
\bottomrule
\multicolumn{3}{@{}l}{\footnotesize \textit{Source: Financial Accounts of the United States (Z.1) and the European System of Accounts (ESA 2010)}}
\end{tabular} 
\end{table}

\section{Financial Instruments}

Financial instruments define the types of claims through which financial institutions provide funding to other sectors. Since the objective of this thesis is to trace financial exposures from financial institutions to real-economy agents, instruments must be classified consistently across the United States and the Eurozone. All instruments are therefore grouped into ten broad categories: Monetary Gold and Special Drawing Rights (SDRs)\footnote{Special Drawing Rights are international reserve assets created by the IMF to supplement the reserves of its member countries.}, Currency and Deposits, Debt Securities, Loans, Equity, Money Market Fund Shares, Mutual Fund Shares, Insurance, Pension, and Other. These categories are denoted by $\alpha \in A$.

As with the classification of financial institutions, the United States reports a more granular breakdown of financial instruments, whereas Eurozone data are generally presented in a more aggregated form. The harmonized classification used in this thesis therefore preserves the main economic function of each instrument while ensuring comparability across the two regions. Detailed descriptions of the instruments included in each category are provided in Table \ref{tab:instrument_mapping}.

\begin{table}[ht!]
\centering
\caption{Financial Instruments: US and Eurozone}
\label{tab:instrument_mapping}
\small
\begin{tabular}{@{} >{\raggedright\arraybackslash}p{3cm} >{\raggedright\arraybackslash}p{6.5cm} >{\raggedright\arraybackslash}p{5.5cm} @{}}
\toprule
\textbf{Instrument Group} & \textbf{United States} & \textbf{Eurozone} \\
\midrule
% ---------------------------------------------------------
% 1. MONETARY GOLD AND SDRs
% ---------------------------------------------------------
\textbf{Monetary Gold and SDRs} & 
Monetary Gold and SDRs & 
Monetary gold and SDRs \\
\midrule
% ---------------------------------------------------------
% 2. CURRENCY AND DEPOSITS
% ---------------------------------------------------------
\textbf{Currency and Deposits} & 
Checkable Deposits and Currency \newline 
Time and Savings Deposits \newline 
Other Deposits \newline 
Net Interbank Transactions & 
Currency \newline
Deposits \\
\midrule
% ---------------------------------------------------------
% 3. DEBT SECURITIES
% ---------------------------------------------------------
\textbf{Debt securities} & 
Treasury Securities \newline 
Agency- and GSE-Backed Securities \newline 
Municipal Securities \newline 
Corporate and Foreign Bonds \newline 
Open Market Paper \newline 
FDI Debt & 
Debt securities \\
\midrule
% ---------------------------------------------------------
% 4. LOANS
% ---------------------------------------------------------
\textbf{Loans} & 
Federal Funds and Repos \newline 
Home Mortgages \newline 
Multifamily Residential Mortgages \newline 
Commercial Mortgages \newline 
Farm Mortgages \newline 
Consumer Credit \newline 
Depository Institution Loans N.E.C. \newline 
Other Loans and Advances & 
Loans \\
\midrule
% ---------------------------------------------------------
% 5. EQUITY
% ---------------------------------------------------------
\textbf{Equity} & 
Corporate Equity \newline 
FDI Equity \newline 
Miscellaneous Other Equity & 
Listed shares \newline 
Unlisted shares \newline 
Other equity \\
\midrule
% ---------------------------------------------------------
% 6. MONEY MARKET FUNDS SHARES
% ---------------------------------------------------------
\textbf{Money Market Funds Shares} & 
Money Market Funds Shares & 
Money market funds shares \\
\midrule
% ---------------------------------------------------------
% 7. MUTUAL FUND SHARES
% ---------------------------------------------------------
\textbf{Mutual Fund Shares} & 
Mutual Fund Shares & 
Non-MMF investment funds \\
\midrule
% ---------------------------------------------------------
% 8. INSURANCE
% ---------------------------------------------------------
\textbf{Insurance} & 
Life Insurance Reserves & 
Life insurance \newline 
Non-life insurance technical reserves and provision for calls under standardised guarantees \\
\midrule
% ---------------------------------------------------------
% 9. PENSION
% ---------------------------------------------------------
\textbf{Pension} & 
Pension Entitlements & 
Pension schemes \\
\midrule
% ---------------------------------------------------------
% 10. OTHER
% ---------------------------------------------------------
\textbf{Other} & 
Trade Credit \newline 
Taxes Payable by Businesses \newline 
Identified Miscellaneous Financial Claims \newline 
Unidentified Miscellaneous Financial Claims & 
Financial derivatives \newline 
Other accounts \\
\bottomrule
\multicolumn{3}{@{}l}{\footnotesize \textit{Source: Financial Accounts of the United States (Z.1) and the European System of Accounts (ESA 2010)}}
\end{tabular} 
\end{table}

\section{From-whom-to-whom matrix}
\label{sec:FWTW_introduction}
Although financial accounts are useful for measuring the size of sectors and instruments, they are not designed to capture the linkages between sectors that provide and those that utilize funding, known as \textit{from-whom-to-whom} relationships. For example, while financial accounts report the total amount of household debt, FWTW data identify which sectors hold these liabilities and the magnitude of their exposures. This information is valuable for understanding financial-system interconnectedness and the accumulation of macro-financial vulnerabilities. In the period leading up to the 2007-2009 financial crisis, for instance, such data could have helped regulators detect earlier the growing share of household debt held outside the traditional banking sector.

This bilateral structure is central to the environmental attribution exercise in this thesis. By identifying which sectors hold claims on which issuers, FWTW matrices make it possible to connect financial institutions as providers of capital to the real-economy sectors whose activities generate emissions and other environmental pressures.



Formally, for each financial instrument $\alpha \in A$, the FWTW matrix takes the form: 
\[
V_t^{(a)}
=
\begin{array}{c}
\text{(Issuers)} \\[2pt]
\left(
\begin{array}{cc}
\begin{array}{c}
V_{FF,t}^{(a)}\\[-3pt]
{\scriptstyle (n_F \times n_F)}
\end{array}
&
\begin{array}{c}
V_{FR,t}^{(a)}\\[-3pt]
{\scriptstyle (n_F \times n_R)}
\end{array}
\\[8pt]
\begin{array}{c}
V_{RF,t}^{(a)}\\[-3pt]
{\scriptstyle (n_R \times n_F)}
\end{array}
&
\begin{array}{c}
V_{RR,t}^{(a)}\\[-3pt]
{\scriptstyle (n_R \times n_R)}
\end{array}
\end{array}
\right)
\end{array}
\hspace{0.3em}
\text{(Holders)}
\]

Rows in $V_t^{(a)}$ correspond to the sectors holding financial claims, while columns correspond to the sectors issuing the corresponding liabilities. The matrix is therefore partitioned according to whether the holder and issuer belong to the financial sector, denoted by $F$, or to the real economy, denoted by $R$. The block $V_{FF,t}^{(a)}$ captures claims among financial institutions, while $V_{FR,t}^{(a)}$ records claims held by financial institutions on real-economy agents. Similarly, $V_{RF,t}^{(a)}$ captures claims held by real-economy agents on financial institutions, and $V_{RR,t}^{(a)}$ records claims within the real economy.

Financial accounts provide the total amount of assets held by a sector (row sums) and the total amount of liability issued by a sector (column sums). Thus, for matrix $V_{t}^{(a)}$, the financial accounts supply: 
\begin{equation}
r^{(a)}_{i,t} = \sum_{j=1}^{N} V^{(a)}_{ij,t}
\; \text{and} \;
c^{(a)}_{j,t} = \sum_{i=1}^{N} V^{(a)}_{ij,t},
\; \text{with} \; N = n_F + n_R
\end{equation}
where $r^{(a)}_{i,t}$ and $c^{(a)}_{j,t}$ denote, respectively, the holdings of sector $i$ and the issuance of sector $j$ for asset $a$. 

To populate the center of the matrix, \textcite{batty_-whom--whom_2023} from the US Federal Reserve starts to fill cells with \textit{"known relationships"}, for which both the holder and issuer are identified. Additionally, they set cells to zero where a particular holder and issuer for that instrument do not directly interact \textit{("exclusion restriction")}. An example is, RoW-RoW cell is always zero because transactions between non-residents will not be captured within domestic financial accounts. Next, they subtract the amounts derived from \textit{"known relationships"} from the total assets and liabilities and assign the remaining holdings proportionally. In other words, they assume that all remaining holders lend to the remaining borrowers in proportion to the remaining overall issuance outstanding for that instrument.      

\textcite{jondeau_environmental_2026} says that the method adopted by the Federal Reserve is similar to the iterative proportional fitting procedure (IPFP), also known as the RAS method (\cite{bacharach_estimating_1965}). The algorithm iteratively adjusts rows and columns as follows:
\begin{equation}
V^{(a)}_{ij,t+1/2}
=
V^{(a)}_{ij,t}
\frac{r^{(a)}_{i,t}}{\sum_{j'=1}^{N} V^{(a)}_{ij',t}}
\quad \text{and} \quad
V^{(a)}_{ij,t+1}
=
V^{(a)}_{ij,t+1/2}
\frac{c^{(a)}_{j,t}}{\sum_{i'=1}^{N} V^{(a)}_{i'j,t+1/2}}
\end{equation}
The procedure can accommodate structural constraints, for example, households neither issue equity nor insurance policies. Under standard conditions, it converges to a unique matrix $V^{(a)*}_{t}$ that exactly matches the observed row and column totals, satisfies the imposed constraints, and remains as close as possible to the initial matrix in sense of maximum entropy.

A key issue in constructing FWTW matrices is that total holdings of an instrument do not always equal total issuance in the financial accounts. Such discrepancies may arise from differences in valuation (e.g., market value for assets versus book value for liabilities), as well as from inconsistencies in data sources, methodologies, or timing \parencite{batty_-whom--whom_2023}. To address this issue, a “discrepancy” sector (DIS) is introduced into the FWTW matrix to reconcile totals with those reported in the financial accounts.

Another issue identified by the Federal Reserve is that FWTW data can occasionally take negative values. This typically occurs when issuance and holdings cannot be separately observed, and only their net difference is recorded as either an asset or a liability in the financial accounts. In such cases, the implied source and use in the FWTW matrix should be interpreted in reverse. Additionally, if the \textit{“known relationships”} for a sector exceed its total issuance or holdings, a negative residual is generated and subsequently distributed across other sectors. In rarer instances, negative values may also reflect data errors. 

